\section{R11：共享存储、TMA swizzle 与固定开销诊断}

\subsection{轮次目标、环境与判定口径}

R10 的 direct-store 结果表明，显式 fragment 到 global 的映射可以做到 exact，
但在完整 non-split K2 上增加了寄存器和地址计算，不能超过默认路径。R11 不再
继续叠加普通 global store，而是拆成三条相互独立的实验线：

\begin{enumerate}
  \item A：删除 direct-output 专用 output ring，检查 union 中的 shared-memory
        预留是否真的能下降；
  \item B：使用 CUTLASS/CuTe 的 128-byte shared-memory swizzle，使 output
        tile 可以由匹配的 SM90 TMA store 写回；
  \item C：用 CUDA event、NCU 和 Nsight Systems 将长序列性能拆成 prepare、
        recurrence、写回以及 CTA/流水线开销，避免把单一 DRAM 指标当成瓶颈结论。
\end{enumerate}

本轮所有正式数据来自同一台 B300（\code{NVIDIA B300 SXM6 AC}，CUDA 12.9，
PyTorch 2.13.0+cu129，SM103a）。benchmark 的固定问题为
\(T=8192,H=96,D=128\)，BF16 state；性能以同一 Python API 和 CUDA events
测量，warmup/iterations/repeats 在日志中保留。R11 的实验开关全部是 build-time
opt-in，默认编译不改变 R10 路径。

\subsection{R11-A：删除 output ring 的实际效果}

\paragraph{设计。}
在 \code{SharedStorageK2} 中，direct-output 路径并不读取 output ring，
因此新增 \code{C1\_K2\_COMPACT\_DIRECT\_STORAGE}，让
\code{OutputStorage::out} 只在需要 shared output pipeline 时保留。该改动不
改变线程映射和 global 地址，只检验资源 union 的下界。开关默认为关闭：

\begin{lstlisting}
FLASH_KDA_C1_DIRECT_OUTPUT=1
FLASH_KDA_C1_DIRECT_OUTPUT_VEC=1
FLASH_KDA_C1_COMPACT_DIRECT_STORAGE=1
\end{lstlisting}

正确性覆盖了 \(T=16,17,64,97\)、\(H=1,4,96\)、BF16/FP32 state、state
in/out、varlen 以及 \(B=2\)。所有 output 和 final state 均逐 bit exact。
在固定 \([8192,96,128]\) 上，compact direct 与同源 control 的 V-split
结果约为 2.017 ms 对 2.017 ms，非 split compact 约 1.67 ms；重复测量不
显示可复现的收益。

NCU launch statistics 仍显示 dynamic shared memory 为 64,640 B（配置值
65,664 B）。原因是 shared-storage union 的峰值由 state/input 或其他成员
主导，删除 output ring 没有改变实际 launch 资源。因此 A 的判定是“代码和
正确性成立，但资源下界不由 output ring 决定”，不接入默认构建。详细结果在
\code{results/r11\_a\_SUMMARY.md} 及其同轮日志中。

\subsection{R11-B：从 standalone swizzle 到完整 K2}

\subsubsection{CuTe layout 的独立验证}

官方示例的关键约束是 TMA descriptor 的 major/global unit-stride mode 必须与
shared tile 的逻辑方向一致。首先编写 \code{r11\_tma\_swizzle\_probe.cu}，
以 \(32\times64\) BF16 tile 测试
\code{Swizzle<3,4,3>}，并显式使用 grid-constant descriptor：

\begin{lstlisting}
template<class TmaStore>
__global__ void kernel(CUTE_GRID_CONSTANT TmaStore const tma_store) { ... }
\end{lstlisting}

未使用 \code{CUTE\_GRID\_CONSTANT} 时，descriptor 在 device 上被错误地按
普通 kernel 参数处理，probe 出现 illegal memory access。修正后 B300 job
24689 输出：

\begin{lstlisting}
logical smem: Sw<3,4,3> ... ((_32,_4),_16):((_1,_512),_32)
tma smem:     Sw<3,4,3> ... (_1,(_32,_4),_16):(_0,(_1,_512),_32)
R11-B swizzled TMA store: mismatches=0 total=2048 cosize=2048 bytes=4096
\end{lstlisting}

这一步证明 128-byte swizzle 与 TMA store 的组合本身可在 B300 exact 运行。

随后用真实 K2 的 \([\text{CHUNK},D]=[16,128]\) 形状直接套用 swizzle。
job 24702 发生 illegal memory access，说明不能仅把原有 row-major 的
\([\text{token},\text{value}]) 逻辑 layout 替换为 swizzle。最终采用与 TMA
major mode 一致的 \([D,\text{CHUNK}]=[128,16]\) shared 逻辑 tile，保留
物理 output pointer 不变，只在 launch 中增加一个 \([H,D,T]\)、stride
\((D,1,DH)\) 的 gmem view。

\subsubsection{完整 K2 的实现}

\code{K2Layouts} 新增 full 和 value-split 两套布局：

\begin{lstlisting}
SwizzledOutputLayout = ComposedLayout<Swizzle<3,4,3>, ...>
SwizzledValueSliceLayout = ComposedLayout<Swizzle<3,4,3>, ...>
TMAFullSwizzledOutputLayout
TMASwizzledValueSliceOutputLayout
\end{lstlisting}

MMA warp 的 16x16 fragment 使用 R9/R10 已验证的逆映射写入 swizzled shared
tile。完整 tile 通过匹配的 TMA store 写回；尾部 \(\text{actual\_len}<16\)
不能缩短 TMA transaction，因此保留 scalar bounded tail。swizzled TMA
descriptor 被加入 kernel 参数，但只有 \code{C1\_K2\_TMA\_SWIZZLED\_OUTPUT=1}
时才实例化实际输出分支。该开关由 \code{setup.py} 转为
\code{-DC1\_K2\_TMA\_SWIZZLED\_OUTPUT=1}，默认环境不设置它。

集成构建命令为：

\begin{lstlisting}
cd <B300_REPO>/assignment02/team/c1_flashkda/FlashKDA
FLASH_KDA_CUDA_ARCHS=103a \\
FLASH_KDA_C1_TMA_SWIZZLED_OUTPUT=1 \\
FLASH_KDA_C1_COMPACT_DIRECT_STORAGE=0 \\
FLASH_KDA_C1_DIRECT_OUTPUT=0 \\
FLASH_KDA_C1_DIRECT_OUTPUT_VEC=0 \\
../../../.venv/bin/python setup.py build_ext --inplace
\end{lstlisting}

首次构建 job 24739 仅因错误使用 \code{../../.venv/bin/python} 失败；使用
\code{../../../.venv/bin/python} 后 job 24746 成功。ptxas 没有报告 spill，
相关 recurrence 实例的 register 数仍在 66--77/thread 范围。

\subsubsection{集成 exactness}

由于 Slurm batch shell 默认没有把 venv 的 ninja 加入 PATH，第一次 exactness
提交只记录了 PyTorch JIT 的环境失败。显式设置
\code{PATH=<assignment02>/.venv/bin:\$PATH} 后，job 24758 完成两组测试：

\begin{table}[htbp]
\centering
\small
\begin{tabular}{l r r r r}
\toprule
测试集 & cases & output diff & state diff & 覆盖范围 \\
\midrule
native smoke & 5 & 0 & 0 & T=16/17/64, H=1/96, state I/O \\
extended smoke & 5 & 0 & 0 & T=97, varlen, B=2, H=4/96, BF16/FP32 \\
\bottomrule
\end{tabular}
\caption{R11-B swizzled TMA 完整 K2 的 exactness。}
\end{table}

这证明 swizzled shared layout、full-tile TMA、tail scalar fallback 和
value-split 分支在已覆盖输入上保持 output/final-state 位级一致；但它不等价
于性能收益。

\subsubsection{同源性能对比}

先以 \code{C1\_K2\_TMA\_SWIZZLED\_OUTPUT=0} 重建 baseline（job 24761），
再以开关为 1 重建 swizzled 版本（job 24768）。两者均清理 build 和旧
\code{flash\_kda\_C*.so}，并使用相同 benchmark：warmup=20、iters=50、
repeats=3。

\begin{table}[htbp]
\centering
\small
\begin{tabular}{l r r r}
\toprule
variant & BF16 state (ms) & no state (ms) & FP32 state (ms) \\
\midrule
same-source baseline & 1.0272 & 1.0285 & 0.9977 \\
swizzled TMA & 1.7623 & 1.7615 & 1.7104 \\
\midrule
relative change & +71.6\% & +71.2\% & +71.4\% \\
\bottomrule
\end{tabular}
\caption{B300，固定 \([8192,96,128]\) 的同源 CUDA-event benchmark。}
\end{table}

为了检查是否只有大 head 数退化，另做 warmup=10、iters=30、repeats=2 的
H=1/4 对照。H=1 的 BF16 state 为 baseline 0.7329 ms、swizzled 0.7332 ms；
H=4 为 0.7479 ms、0.7480 ms。小 H 几乎没有变化，说明这些实例没有触发
完整的高并行 swizzled output 代价；H=96 才暴露出该路径的主要问题。

\begin{decision}{R11-B 判定}
128-byte swizzle + TMA store 在独立 probe 和完整 K2 上均可 exact，但当前
实现将 output fragment 额外 materialize 到 shared tile，再等待 TMA store
完成；对 H=96 增加约 72\% 端到端延迟。它不能替代默认路径，保留为显式
\code{C1\_K2\_TMA\_SWIZZLED\_OUTPUT} 负结果/后续优化基线。下一次若重访，
必须消除额外 materialization 或把它与 MMA/producer pipeline 融合，而不是
继续调整 swizzle 参数。
\end{decision}

\subsection{R11-C：固定开销与 CTA 形态诊断}

\subsubsection{测量方法}

\code{r11\_c\_overhead\_probe.py} 使用 public \code{flash\_kda.fwd} API，
CUDA event 包住一次 prepare+recurrence；CPU allocation 与 Python dispatch
不计入 GPU event。NCU basic set 收集 grid/block、register、dynamic shared
memory、active warps 和 DRAM bytes；Nsight Systems 记录短序列和长序列的
kernel 时间。NCU replay 时间只用于结构观察，不参与正式性能排序。

\subsubsection{H=96 的 V-split 结构}

\begin{table}[htbp]
\centering
\small
\begin{tabular}{l c c c c}
\toprule
variant & grid & block & dynamic smem & registers/thread \\
\midrule
baseline & (1,96,1) & 192 & 98,432 B & 66--74 \\
V-split & (1,192,1) & 128 & 68,608 B & 78--82 \\
\bottomrule
\end{tabular}
\caption{R11-C NCU launch statistics，H=96。}
\end{table}

V-split 把 CTA 数量翻倍、每个 CTA 的计算 warp 数减少；虽然 dynamic shared
memory 下降约 30\%，其 recurrence kernel 反而从约 1.24--1.29 ms 增至
1.56--1.58 ms。完整 API 的 T=8192 CUDA-event 结果如下：

\begin{table}[htbp]
\centering
\small
\begin{tabular}{l r r r}
\toprule
H & baseline (ms) & V-split (ms) & V-split change \\
\midrule
1 & 1.2680 & 1.1094 & -12.5\% \\
4 & 1.2916 & 1.1518 & -10.8\% \\
96 & 1.7582 & 2.0066 & +14.1\% \\
\bottomrule
\end{tabular}
\caption{R11-C 同一 B300 环境的正式 CUDA-event 对照。}
\end{table}

H=96 的 NCU DRAM throughput 在 V-split 中更高，但端到端更慢，因而不能把
“DRAM read/write 较少”直接解释为总延迟瓶颈已经被解决。Nsight Systems 还
显示 T=16 时总 GPU 时间只有约 30--60\(\mu\)s，固定 launch、barrier 和
pipeline 成本占比显著；varlen 的加速主要来自多个 sequence CTA 并行，关键
路径接近最长 sequence，而不是简单按总 token 数缩放。

\subsection{R11 总结与下一轮边界}

本轮得到三个可复核结论：

\begin{enumerate}
  \item 删除 direct-output ring 没有降低实际 dynamic shared memory，说明
        union 峰值由其他 storage 主导；
  \item swizzled shared/TMA store 已从独立 probe 迁移到真实 K2 并通过完整
        exactness，但 H=96 慢约 72\%，因此只能保留为 opt-in 负结果；
  \item V-split 的资源节省与端到端收益不单调，H=1/4 有约 11--13\% 收益，
        H=96 却慢约 14\%，主要矛盾是 CTA/warp 形态而非单一 DRAM 计数器。
\end{enumerate}

R11 不改变默认构建：\code{C1\_K2\_TMA\_SWIZZLED\_OUTPUT}、
\code{C1\_K2\_COMPACT\_DIRECT\_STORAGE} 和 R10 direct 开关均保持关闭。
后续优化若继续，应优先研究 H=96 的 CTA 融合/warp-specialized epilogue 与
真实 output-storage 生命周期；不应把当前 swizzled materialization 直接
接入默认路径，也不应仅凭 shared-memory 或 DRAM 指标宣称已达到题目目标。

